\chapter*{Introduction}
\addcontentsline{toc}{chapter}{Introduction}
\markboth{Introduction}{}

Le secteur extractif occupe une place croissante dans l'économie sénégalaise. Au-delà des phosphates et, plus récemment, des hydrocarbures du bassin sénégalo-mauritanien, l'or constitue depuis le début des années 2010 le principal produit minier exporté par le pays, porté essentiellement par le complexe de Sabodala, dans la région de Kédougou, à l'extrême sud-est du territoire. Membre de l'Initiative pour la Transparence dans les Industries Extractives (ITIE) depuis 2013 et reconnu conforme aux exigences de la Norme ITIE, le Sénégal publie chaque année des rapports de conciliation qui détaillent les paiements effectués par les sociétés extractives et les recettes perçues par l'État. Ces rapports renseignent sur les flux \emph{constatés}, mais ne permettent pas, à eux seuls, d'apprécier si le régime fiscal applicable à un projet donné permet à l'État de capter une part « raisonnable » de la rente minière, ni de situer ce régime par rapport aux pratiques observées dans d'autres juridictions minières.

C'est précisément l'objet des modèles de simulation fiscale de type \emph{Discounted Cash Flow} (DCF), au premier rang desquels figure le modèle \emph{Fiscal Analysis of Resource Industries} (FARI) développé par le Département des finances publiques du Fonds Monétaire International. Diffusé sous une forme simplifiée par le \emph{Natural Resource Governance Institute} (NRGI), ce type d'outil permet de simuler, sur toute la durée de vie d'un projet, la répartition des flux de trésorerie entre l'investisseur et l'État, ainsi que des indicateurs de synthèse tels que la Valeur Actuelle Nette (\VAN), le Taux de Rentabilité Interne (\TRI) et le Taux Effectif Moyen d'Imposition (\TEMI). Le mémoire de \citet{kourouma2019}, réalisé à l'Université du Québec à Montréal dans le cadre d'un stage au sein du projet de Gouvernance Régionale du Secteur Extractif en Afrique de l'Ouest (GRSE) de la GIZ, illustre cette démarche en l'appliquant au projet bauxitique guinéen d'ALUFER pour comparer les codes miniers guinéens de 1995 et de 2011 amendé à un régime fiscal hypothétique.

Le présent rapport reprend cette même famille d'outils --- en l'occurrence le modèle financier d'une mine d'or diffusé par NRGI, dont l'architecture suit la méthodologie FARI --- et l'applique pour la première fois, à notre connaissance, au cas sénégalais, en utilisant comme support empirique le projet aurifère Sabodala-Massawa exploité par Sabodala Gold Operations (SGO). Les données opérationnelles et financières du projet sont extraites de l'étude de préfaisabilité (\emph{Pre-Feasibility Study}, PFS) publiée en août 2020 par Lycopodium pour le compte de Teranga Gold Corporation (aujourd'hui intégrée à Endeavour Mining), qui détaille notamment le calendrier de production sur la période 2020-2036, les dépenses d'investissement de développement et de maintien, les coûts opérationnels et le régime fiscal applicable au projet.

\paragraph{Question de recherche.} Dans quelle mesure le régime fiscal effectivement appliqué au projet Sabodala-Massawa permet-il à l'État sénégalais de capter une part de la rente minière comparable à celle observée dans d'autres juridictions aurifères, et quelle est la structure de ces recettes par instrument fiscal~?

\paragraph{Démarche.} Pour répondre à cette question, le modèle NRGI --- qui intègre par défaut une comparaison entre le code minier de la République Démocratique du Congo (RDC) et neuf autres régimes de référence (Chili, Pérou, Australie-Occidentale, Tanzanie, Ghana, Zambie, ainsi que trois variantes du régime congolais) --- a été complété par un onzième scénario reproduisant les paramètres fiscaux applicables à SGO~: redevance de 5\,\% sur la valeur des expéditions d'or, impôt sur les sociétés de 25\,\%, participation gratuite de l'État de 10\,\% et exonération de droits de douane liée au statut d'entreprise franche d'exportation. Le profil opérationnel générique initialement présent dans l'outil (une mine fictive de taille moyenne) a quant à lui été remplacé par le profil réel de Sabodala-Massawa sur la période 2020-2036, reconstitué à partir des tableaux de synthèse du PFS.

\paragraph{Organisation du rapport.} Le chapitre~1 situe le cadre théorique de l'analyse en présentant les enjeux de la fiscalité minière dans les pays riches en ressources, la typologie des instruments fiscaux mobilisables et les principes du modèle FARI. Le chapitre~2 décrit le cadre légal et fiscal du secteur aurifère sénégalais et présente le projet Sabodala-Massawa. Le chapitre~3 détaille la méthodologie suivie pour adapter le modèle NRGI au cas sénégalais, en particulier l'extraction des données du PFS et la construction du scénario fiscal « Sénégal ». Le chapitre~4 présente et discute les résultats obtenus --- indicateurs de rentabilité, structure des recettes publiques et positionnement du \TEMI sénégalais dans la comparaison internationale --- ainsi que les limites méthodologiques rencontrées. Le chapitre~5 propose des pistes pour l'utilisation de ce type d'outil par l'ITIE Sénégal. Le rapport se termine par une conclusion générale et des annexes présentant le détail des données utilisées.
